\documentclass[11pt,letterpaper]{article}
\usepackage{technical-implementation}
% Shared non-semantic helper definitions; one semantic module is loaded by each root.
\newcommand{\PTSource}[1]{\par\smallskip{\footnotesize\color{Meta}\textbf{Book crosswalk:} #1}\par\smallskip}
\newcommand{\PTOutcome}[1]{\begin{tipbox}[Competency outcome]#1\end{tipbox}}
\newcommand{\PTEvidence}[1]{\begin{notebox}[Required evidence]#1\end{notebox}}
\newcount\ptlinecount
\newcommand{\PTLines}[1]{\par\ptlinecount=0\loop\ifnum\ptlinecount<#1\par\noindent\rule{\linewidth}{0.25pt}\vspace{0.22cm}\advance\ptlinecount by1\repeat\par}
\newcommand{\PTField}[1]{\noindent\textbf{#1}\quad\rule{0.66\linewidth}{0.3pt}\par\vspace{0.22cm}}
\newcommand{\PTRating}{\noindent\textbf{Rating:}\quad 0\quad 1\quad 2\quad 3\quad 4\qquad\textbf{Evidence ref:}\rule{0.30\linewidth}{0.3pt}\par\vspace{0.14cm}}
\newcommand{\PTTask}[2]{\par\medskip\noindent\textbf{Task #1.} #2\par}

\setDocTitle{Chapter 15: Non-Chomsky Grammars and Their Parsers}
\setDocSubtitle{Foundations, reasoning, guided experiments and assessment}
\setDocVersion{1.0}\setDocStatus{Engineering education / source-aligned}
\setDocOwner{Compiler Engineering}
\setDocDate{30 October 2026}
\setDocSummary{Chapter 15 learning companion.}
\begin{document}\makedoctitle\makedocsummary
\tableofcontents\clearpage
\section{Chapter orientation}
This guide follows Chapter 15 of Grune and Jacobs, \emph{Parsing Techniques: A Practical Guide}, second edition. Printed pages 473--520. The focus is on reconstructing the chapter's reasoning, verifying explicit assumptions, and translating its central idea into a controlled engineering experiment. This is an independently authored learning guide, not a substitute for the book.
\PTSource{Chapter 15, Non-Chomsky Grammars and Their Parsers; printed pages 473--520.}
\section{Learning objectives}
\begin{enumerate}
\item Define and differentiate two-level grammars using a small grammar and its accepted language.
\item Define and differentiate attribute and affix grammars using a small grammar and its accepted language.
\item Construct an example or counterexample illustrating tree-adjoining grammars.
\item Evaluate the operational constraints of ordered, recognition-system, and boolean grammars and retain reproducible evidence.
\end{enumerate}
\section{Core topics and review questions}
\subsection{Two-level grammars}
Write the exact input domain, grammar assumptions, data structures and observable outputs. Distinguish a mathematical language property from the behavior of one implementation.
\textbf{Self-test:} Which invariant or decision criterion makes this topic precise, and what input would falsify an incorrect implementation?
\subsection{Attribute and affix grammars}
Construct two contrasting examples. Give a complete trace of states or sentential forms and indicate the condition under which the process halts.
\textbf{Self-test:} Which invariant or decision criterion makes this topic precise, and what input would falsify an incorrect implementation?
\subsection{Tree-adjoining grammars}
Identify a counterexample, ambiguity, cycle or boundary condition; explain whether it is a grammar property, an algorithm limitation, or an implementation bug.
\textbf{Self-test:} Which invariant or decision criterion makes this topic precise, and what input would falsify an incorrect implementation?
\subsection{Ordered, recognition-system, and Boolean grammars}
State the cost model, failure modes and evidence needed for a defensible design choice. Do not extrapolate results beyond the grammar family studied.
\textbf{Self-test:} Which invariant or decision criterion makes this topic precise, and what input would falsify an incorrect implementation?
\section{Technical deep dive: syntax and attributes}
An attribute grammar associates computed properties with syntax nodes. Synthesized attributes flow from children to parents, while inherited attributes can transport contextual information downward or sideways under a dependency order. For an arithmetic AST, a synthesized value can evaluate $E\to E+T$ from child values; for name binding, an inherited symbol environment is often clearer. Analyze cycles in attribute dependencies independently from syntactic recursion.
\section{Worked reasoning prompt}
Compute a synthesized attribute for evaluating a simple arithmetic expression tree. Derive each intermediate state by hand; write the invariant that justifies advancing to the next state. Identify the accepted result and one structurally similar rejected result.
\section{Implementation and laboratory assignment}
Create a comparison matrix for at least four expressive extensions and their parsing implications. Record the grammar, input corpus, expected outputs, maximum input size and a regression test for the principal failure mode. For each test, independently justify the expected result rather than reading it from the program.
\section{Failure analysis and technical discussion}
Do not assume that semantic constraints should be folded into unrestricted context-sensitive grammar rules. Explain how your implementation would detect the failure and which diagnostics can distinguish grammar invalidity from resource exhaustion.
\section{Evidence checklist}
\begin{itemize}
\item Grammar and accepted-language specification with version and fixtures.
\item Hand-executed trace and implementation trace for the same accepted and rejected cases.
\item At least one adversarial or boundary input and an explanation of the expected behavior.
\item Correctness argument, complexity claim with assumptions, and unit-test output.
\item Cross-reference to a related chapter and one explicit unresolved question.
\end{itemize}
\section{Independent examination}
\begin{enumerate}
\item Explain the algorithm without using implementation-language details.
\item Give a proof sketch or invariant for acceptance and termination on the specified grammar class.
\item Change one grammar rule and predict which tests must change before running them.
\item Compare behavior on malformed input, ambiguous input, and input beyond configured resource limits.
\end{enumerate}

\end{document}
